% =============================================================================
% CENG 467 – Natural Language Understanding and Generation
% Take-Home Midterm Examination
% Student ID: 310201050
% Instructor: Prof. Dr. Aytuğ Onan
% Date: April 2026
% =============================================================================

\documentclass[11pt,a4paper]{article}

% ── Layout & encoding ────────────────────────────────────────────────────────
\usepackage[T1]{fontenc}
\usepackage[utf8]{inputenc}
\usepackage[english]{babel}
\usepackage[a4paper, margin=2.5cm]{geometry}
\usepackage{microtype}
\usepackage{setspace}
\onehalfspacing

% ── Mathematics ──────────────────────────────────────────────────────────────
\usepackage{amsmath, amssymb, amsthm}

% ── Tables / figures ─────────────────────────────────────────────────────────
\usepackage{booktabs}
\usepackage{graphicx}
\usepackage{caption}
\usepackage{subcaption}
\usepackage{array}

% ── Code listings (optional) ─────────────────────────────────────────────────
\usepackage{listings}
\usepackage{xcolor}
\definecolor{codegray}{gray}{0.95}
\lstset{
    basicstyle      = \ttfamily\small,
    backgroundcolor = \color{codegray},
    breaklines      = true,
    frame           = single,
    numbers         = left,
    numberstyle     = \tiny,
    keywordstyle    = \color{blue},
    commentstyle    = \color{gray},
    stringstyle     = \color{red!70!black},
    showstringspaces= false,
}

% ── Hyperlinks ───────────────────────────────────────────────────────────────
\usepackage[colorlinks=true, linkcolor=black, citecolor=blue, urlcolor=blue]{hyperref}
\usepackage{url}

% ── Bibliography ─────────────────────────────────────────────────────────────
\usepackage[numbers, sort&compress]{natbib}

% =============================================================================
% TITLE PAGE
% =============================================================================
\title{
    \vspace{-1cm}
    \large CENG 467 — Natural Language Understanding and Generation \\
    \vspace{0.4cm}
    \LARGE\textbf{Take-Home Midterm Examination} \\
    \vspace{0.6cm}
    \large A Comparative Study of Representation, Sequence Labeling, \\
    Summarization, Translation and Language Modeling
}
\author{
    \textbf{Student ID:} 310201050 \\
    \vspace{0.2cm}
    \small Instructor: Prof.\ Dr.\ Aytu\u{g} Onan \\
    \small \.{I}zmir Institute of Technology --- Department of Computer Engineering
}
\date{April 2026}

% =============================================================================
\begin{document}
\maketitle

\thispagestyle{empty}
\vspace{0.5cm}

\begin{abstract}
\noindent
This report presents a controlled comparative study of five core
NLP tasks: text classification, named entity recognition,
text summarization, machine translation, and language modeling.
For each task, two or three model families are evaluated on a
shared benchmark dataset under fixed random seeds, consistent
preprocessing pipelines, and held-out test sets that are used only
once for final evaluation. Sparse, dense, and contextual
representations are compared head-to-head, and qualitative error
analyses are provided alongside the quantitative results. Across
all five tasks, contextual representations from pretrained
Transformer models outperform classical and from-scratch recurrent
baselines, with the F1 gap reaching nearly 20 points for named
entity recognition and a 10$\times$ perplexity reduction for
language modelling. All experiments are reproducible from the
accompanying code repository on an Apple M2 (8\,GB) device with
seed 3102.
\end{abstract}

\vspace{0.5cm}
\noindent
\textbf{Repository:}
\url{https://github.com/<your-github-username>/CENG467_Midterm_310201050}

\vspace{1cm}
\tableofcontents

\newpage

% =============================================================================
% INTRODUCTION
% =============================================================================
\section{Introduction}

This document presents the solutions to the five questions of the
CENG 467 take-home midterm. Each section is self-contained, including
its own dataset description, methodology, hyperparameters, results,
and discussion.

The study is conducted under three overarching constraints. First,
all experiments use a fixed random seed (3102, derived from the
author's student identifier) to ensure exact reproducibility.
Second, dataset splits are kept consistent between models within a
single question; the test set is touched exactly once at the end of
each pipeline. Third, all training and inference is performed on a
single Apple~M2 device with 8\,GB of RAM, which informs the choice
of distilled and small-footprint models throughout the study.

The complete source code is organised by question
(\texttt{q1\_classification/}, \texttt{q2\_ner/}, etc.) and is
available in the accompanying GitHub repository. Reproduction
instructions are provided in the \texttt{README.md} file.

\subsection{Headline Results}

Table~\ref{tab:headline} summarises the principal finding from each
question. In every case, the contextual or pretrained Transformer
model outperforms the from-scratch baseline; the gap is largest where
training data is most limited (NER) and on tasks where surface
features are weakest (language modelling).

\begin{table}[h]
\centering
\caption{Best result on each task and the from-scratch baseline.}
\label{tab:headline}
\small
\begin{tabular}{llcc}
\toprule
Task & Headline metric & Baseline & Best (pretrained) \\
\midrule
Q1 Classification (IMDb)        & Macro-F1   & 0.7488 (BiLSTM)       & \textbf{0.8970} (DistilBERT) \\
Q2 NER (CoNLL-2003)             & F1         & 0.7011 (BiLSTM-CRF)   & \textbf{0.8975} (DistilBERT) \\
Q3 Summarization (CNN/DM)       & ROUGE-L    & 0.1883 (LexRank)      & \textbf{0.2580} (DistilBART) \\
Q4 Translation (Multi30k de–en) & BLEU       & 27.69 (Seq2Seq)       & \textbf{39.66} (MarianMT)    \\
Q5 Language Modeling (WikiText-2) & Test PPL & 3587.97 (Trigram)     & \textbf{353.41} (LSTM)       \\
\bottomrule
\end{tabular}
\end{table}

% =============================================================================
% INCLUDE EACH QUESTION
% =============================================================================
% ─────────────────────────────────────────────────────────────────
% CENG 467 Midterm – Question 1: Text Classification
% Student ID: 310201050
% ─────────────────────────────────────────────────────────────────

\section{Question 1: Representation Learning in Text Classification}

\subsection{Dataset}

The IMDb Large Movie Review Dataset \cite{maas2011imdb} consists of 50{,}000
binary-labelled movie reviews (25{,}000 train / 25{,}000 test), evenly split
between positive and negative sentiment classes. Because fine-tuning a
Transformer on the full training corpus exceeds the available memory budget
of the Apple~M2 (8\,GB RAM), a reproducible subset was drawn using seed 3102
derived from the author's student identifier: $6{,}000$ training samples,
$1{,}000$ validation samples, and $1{,}000$ held-out test samples. The test
set was used \emph{exactly once} for final evaluation.

\subsection{Preprocessing and Tokenisation Strategies}

A shared normalisation pipeline is applied to all texts before
model-specific tokenisation. HTML tags are stripped with a regular
expression, the text is lower-cased, all non-alphabetic characters are
removed, and consecutive whitespace is collapsed to a single space. Three
tokenisation strategies are then applied, each suited to the
representational assumptions of a different model.

\paragraph{Strategy A — Whitespace Split with Stopword Removal.}
Tokens are produced by splitting on whitespace after normalisation,
followed by removal of English stopwords (NLTK list, 179 tokens). This
yields compact, content-bearing token sequences and is used by the TF-IDF
representation. Removing stopwords reduces feature noise but discards
negation words such as \textit{not}; the TF-IDF bigram range partially
compensates for this loss.

\paragraph{Strategy B — NLTK \texttt{word\_tokenize}.}
NLTK's rule-based tokeniser is applied after normalisation \emph{without}
stopword removal, preserving punctuation-adjacent morphology and
contractions that carry affective weight (e.g.\ \textit{wasn't},
\textit{couldn't}). This strategy is used by BiLSTM to retain richer
sequential cues.

\paragraph{Strategy C — WordPiece subword tokenisation.}
DistilBERT employs its own WordPiece subword tokeniser via
\texttt{DistilBertTokenizerFast}, which decomposes out-of-vocabulary words
into known sub-units, making it the most robust of the three approaches.

\subsection{Models}

\subsubsection{TF-IDF + Logistic Regression}

A sparse TF-IDF matrix is constructed from Strategy-A tokens with $30{,}000$
maximum features, unigram and bigram ranges, and sublinear TF scaling
($\text{tf} \leftarrow 1+\log\text{tf}$). A Logistic Regression classifier
with $C=0.7$ is trained via L-BFGS. This constitutes the \emph{sparse}
representation baseline.

\subsubsection{BiLSTM}

A two-layer bidirectional LSTM \cite{hochreiter1997lstm} is implemented with
the configuration shown in Table~\ref{tab:bilstm_hp}. The final hidden
states from the last forward and backward layers are concatenated and
passed through a linear classification head. This constitutes the
\emph{dense sequential} representation.

\begin{table}[h]
\centering
\caption{BiLSTM Hyperparameters}
\label{tab:bilstm_hp}
\begin{tabular}{ll}
\toprule
Hyperparameter & Value \\
\midrule
Vocabulary size     & 20{,}000 \\
Embedding dimension & 128 \\
Hidden dimension    & 128 \\
Number of layers    & 2 (bidirectional) \\
Dropout             & 0.35 \\
Maximum sequence length & 256 \\
Batch size          & 32 \\
Epochs              & 5 \\
Optimiser / LR      & Adam / $10^{-3}$ \\
Gradient clipping   & 1.0 \\
\bottomrule
\end{tabular}
\end{table}

\subsubsection{DistilBERT}

DistilBERT-base-uncased \cite{sanh2019distilbert} is fine-tuned for sequence
classification (Table~\ref{tab:bert_hp}). The full self-attention stack
provides \emph{contextual} representations where each token's embedding is
conditioned on the entire input.

\begin{table}[h]
\centering
\caption{DistilBERT Fine-Tuning Hyperparameters}
\label{tab:bert_hp}
\begin{tabular}{ll}
\toprule
Hyperparameter & Value \\
\midrule
Pretrained model        & \texttt{distilbert-base-uncased} \\
Maximum sequence length & 256 \\
Batch size              & 16 \\
Epochs                  & 3 \\
Learning rate           & $2 \times 10^{-5}$ \\
Weight decay            & 0.01 \\
Warmup ratio            & 10\% \\
Gradient clipping       & 1.0 \\
Scheduler               & Linear with warmup \\
\bottomrule
\end{tabular}
\end{table}

\subsection{Evaluation Results}

All models are evaluated on the same held-out test split ($n=1{,}000$) using
Accuracy and Macro-F1. Given the balanced class distribution, both metrics
yield consistent estimates of generalisation
(Table~\ref{tab:q1_results}).

\begin{table}[h]
\centering
\caption{Test-Set Performance on IMDb (seed 3102)}
\label{tab:q1_results}
\begin{tabular}{lcc}
\toprule
Model & Accuracy & Macro-F1 \\
\midrule
TF-IDF + Logistic Regression & 0.8550 & 0.8548 \\
BiLSTM                       & 0.7490 & 0.7488 \\
DistilBERT                   & \textbf{0.8970} & \textbf{0.8970} \\
\bottomrule
\end{tabular}
\end{table}

The pattern is striking: the sparse TF-IDF baseline (85.5\%) substantially
outperforms the dense BiLSTM (74.9\%), while the contextual DistilBERT
(89.7\%) outperforms both. The strong showing of TF-IDF reflects the
fact that IMDb reviews are long and lexically rich in sentiment-bearing
unigrams, exactly the regime in which sparse linear models excel. The
BiLSTM, despite its sequential modelling capacity, is data-hungry: with
only 6{,}000 randomly initialised embeddings, it cannot generalise as
well as a well-regularised linear model on the same data.

\subsection{Misclassification Analysis}

Five misclassified examples produced by DistilBERT were inspected in
detail; the full text of each is preserved in
\texttt{q1\_results.json}. Three of the five errors share a single
recurring pattern.

\paragraph{Pattern 1 — Ironic praise of niche or ``cult'' material.}
Reviews [8], [14], and [31] are all positive reviews predicted as
negative. They share a rhetorical strategy in which the reviewer
enthusiastically endorses a film while using ironic or self-deprecating
descriptors (\textit{``Citizen Cane this film is not''}, \textit{``a
mess''}, \textit{``bad hair''}, \textit{``about as good as a bad
dino-flick can get''}). DistilBERT picks up on the surface-level negative
adjectives and misclassifies the overall stance, which is unmistakably
positive on a careful reading.

\paragraph{Pattern 2 — Mixed sentiment with a reversed conclusion.}
Review [50] is a negative review predicted as positive: the writer opens
by criticising the film as boring and predictable, then spends most of
the body praising specific actors and on-screen chemistry, before
returning to the negative verdict. DistilBERT's 256-token window captures
predominantly the central praise, missing the bracketing critical frame.

\paragraph{Pattern 3 — Affective negation with mixed lexical signal.}
Review [1] explicitly disclaims itself (\textit{``this film will never
win an oscar''}, \textit{``Citizen Cane this film is not''}) but
expresses genuine enjoyment (\textit{``I had fun''}, \textit{``I liked
the cheese''}). The lexical signal is genuinely conflicted, and the
model resolves it incorrectly.

These three patterns —- irony, mixed-sentiment with reversal, and
affective negation —- are all variants of a single underlying problem:
\emph{the model is reading lexical polarity rather than rhetorical
stance}. Such errors are well-documented for fine-tuned Transformers and
are not eliminable without explicit modelling of discourse structure or
substantially larger training data.

\subsection{Discussion: Representation Types}

\paragraph{TF-IDF (Sparse).} Despite its simplicity, TF-IDF with logistic
regression delivered $85.5\%$ accuracy --- well within ten points of
the DistilBERT result --- because IMDb reviews are long and
sentiment-rich at the unigram level. The model is computationally cheap,
fully interpretable through feature weights, and requires no GPU.
However, it is completely insensitive to word order and context; surface
forms such as \textit{good} and \textit{not good} receive independent
representations.

\paragraph{BiLSTM (Dense Sequential).} Surprisingly, the BiLSTM
\emph{underperforms} TF-IDF on this subset by more than 10 points. This
counterintuitive ordering reflects two factors. First, the BiLSTM is
initialised with random embeddings, so it must learn its representations
from only 6{,}000 examples --- insufficient to outperform a
well-regularised sparse linear model. Second, the long-tail vocabulary
of IMDb reviews is partially clipped to 20{,}000 types, and rare but
sentiment-bearing words are mapped to \texttt{<UNK>}. With pre-trained
GloVe or word2vec initialisation the result would likely improve, but
even then the gap to DistilBERT would persist.

\paragraph{DistilBERT (Contextual).} Pre-trained self-attention provides
rich, context-dependent representations that markedly outperform both
sparse and recurrent baselines. Each token is encoded with global
context at no additional training cost beyond fine-tuning. The
trade-off is computational: DistilBERT requires substantially more
memory and training time than the other two approaches. On the M2 the
3-epoch fine-tune took roughly $45$\,minutes, compared with seconds for
the sparse model.
\newpage
% ─────────────────────────────────────────────────────────────────
% CENG 467 Midterm – Question 2: Named Entity Recognition
% Student ID: 310201050
% ─────────────────────────────────────────────────────────────────

\section{Question 2: Named Entity Recognition}

\subsection{Dataset and Annotation Structure}

The CoNLL-2003 English shared-task corpus \cite{sang2003conll} is used. It
consists of newswire articles from the Reuters news collection annotated
with four entity types: Persons (\textsc{PER}), Organisations
(\textsc{ORG}), Locations (\textsc{LOC}) and Miscellaneous entities
(\textsc{MISC}). Annotations follow the BIO scheme (\textsc{B}-prefix
for the first token of an entity, \textsc{I}-prefix for inside tokens, and
\textsc{O} for non-entity tokens), yielding a 9-class tag set:
$\{$O, B-PER, I-PER, B-ORG, I-ORG, B-LOC, I-LOC, B-MISC, I-MISC$\}$.

The corpus is loaded from a static mirror because the HuggingFace
\texttt{datasets} library (v3.0+) no longer supports script-based
loaders and all official mirrors of CoNLL-2003 are script-based. The
downloaded files are normalised to BIO tagging by an idempotent
IOB1$\rightarrow$BIO conversion routine. The standard splits are then
used as released: 14{,}041 training sentences, 3{,}250 validation
sentences, and 3{,}453 test sentences. The test set is used
\emph{exactly once} for final evaluation.

\subsection{Token--Label Alignment}

For the BiLSTM-CRF the alignment is trivial because the corpus is already
pre-tokenised at the word level. For DistilBERT, however, each
pre-tokenised word is further decomposed into one or more WordPiece
sub-tokens. We adopt the canonical alignment scheme: \emph{the first
sub-token of each word inherits the original BIO label, while subsequent
sub-tokens are assigned the special ignore index $-100$}. The
\textsc{CrossEntropyLoss} skips these positions, and during evaluation
only the first sub-token of each word contributes to the predicted
sequence, ensuring that entity-level metrics are computed against
word-aligned predictions.

\subsection{Models}

\subsubsection{BiLSTM-CRF (Hybrid)}

A single-layer bidirectional LSTM with a Conditional Random Field
decoding layer \cite{lample2016neural} is trained from scratch. The
CRF layer explicitly models tag-transition probabilities, which
discourages illegal sequences such as \texttt{O B-PER I-LOC}.

\begin{table}[h]
\centering
\caption{BiLSTM-CRF Hyperparameters}
\label{tab:bilstm_crf_hp}
\begin{tabular}{ll}
\toprule
Hyperparameter & Value \\
\midrule
Vocabulary size           & $10{,}951$ words (min freq.\ 2) \\
Embedding dimension       & 96 \\
Hidden dimension          & 192 \\
Number of LSTM layers     & 1 (bidirectional) \\
Dropout                   & 0.40 \\
Batch size                & 32 \\
Epochs                    & 6 \\
Optimiser / LR            & Adam / $8 \times 10^{-4}$ \\
Weight decay              & $10^{-5}$ \\
Gradient clipping         & 1.0 \\
\bottomrule
\end{tabular}
\end{table}

\subsubsection{DistilBERT (Contextual)}

\texttt{distilbert-base-cased} is fine-tuned for token classification.
The \emph{cased} variant is preferred because capitalisation is a strong
NER feature in English news text (\textit{Apple} vs.\ \textit{apple}).

\begin{table}[h]
\centering
\caption{DistilBERT Fine-Tuning Hyperparameters}
\label{tab:bert_ner_hp}
\begin{tabular}{ll}
\toprule
Hyperparameter & Value \\
\midrule
Pretrained model        & \texttt{distilbert-base-cased} \\
Maximum sequence length & 128 \\
Batch size              & 16 \\
Epochs                  & 3 \\
Learning rate           & $3 \times 10^{-5}$ \\
Weight decay            & 0.01 \\
Warmup ratio            & 10\% \\
Gradient clipping       & 1.0 \\
Scheduler               & Linear with warmup \\
\bottomrule
\end{tabular}
\end{table}

\subsection{Evaluation Results}

Entity-level Precision, Recall and F1 are computed using the
\texttt{seqeval} library, which evaluates \emph{exact span matches}: a
prediction is correct only if both the boundaries and the entity type
match the gold annotation (Table~\ref{tab:q2_results}).

\begin{table}[h]
\centering
\caption{CoNLL-2003 Test-Set Results (entity-level, seed 3102)}
\label{tab:q2_results}
\begin{tabular}{lccc}
\toprule
Model & Precision & Recall & F1 \\
\midrule
BiLSTM-CRF         & 0.8010 & 0.6234 & 0.7011 \\
DistilBERT (cased) & \textbf{0.8895} & \textbf{0.9056} & \textbf{0.8975} \\
\bottomrule
\end{tabular}
\end{table}

The contextual DistilBERT outperforms the from-scratch BiLSTM-CRF by
nearly 20 F1 points, the largest margin observed across all five
questions in this study. A per-class breakdown is given in
Table~\ref{tab:q2_per_class}.

\begin{table}[h]
\centering
\caption{Per-Entity F1 Scores on the Test Set}
\label{tab:q2_per_class}
\begin{tabular}{lcc}
\toprule
Entity Class & BiLSTM-CRF F1 & DistilBERT F1 \\
\midrule
\textsc{PER}  & 0.7258 & \textbf{0.9507} \\
\textsc{LOC}  & 0.7777 & \textbf{0.9245} \\
\textsc{ORG}  & 0.6258 & \textbf{0.8754} \\
\textsc{MISC} & 0.6311 & \textbf{0.7682} \\
\bottomrule
\end{tabular}
\end{table}

The relative gap between models is largest for \textsc{ORG} and
\textsc{PER}, which contain the highest density of rare proper nouns,
and smallest for \textsc{MISC}, which is intrinsically the noisiest
entity class.

\subsection{Error Analysis}

The 5{,}644 gold entities in the test set produced the following
DistilBERT error breakdown:

\begin{table}[h]
\centering
\caption{DistilBERT Error Categories on the Test Set}
\label{tab:q2_errors}
\begin{tabular}{lr}
\toprule
Category & Count \\
\midrule
True Positives  & 5{,}106 \\
Type confusion  & 250 \\
False positives & 236 \\
False negatives & 217 \\
Boundary errors & 71 \\
\bottomrule
\end{tabular}
\end{table}

Three observations follow from this distribution.

\paragraph{Boundary errors are remarkably rare.}
Only 71 of the 5{,}644 gold spans (1.3\%) suffered from boundary
mismatches, all but one of which involved \textsc{MISC} entities (e.g.\
predicting only the head of a multi-word entity). The CRF layer of the
BiLSTM-CRF model was originally introduced to suppress exactly this
class of error, but DistilBERT's softmax classifier achieves a
comparable result through context alone.

\paragraph{Type confusion is the dominant error mode.}
With 250 cases (4.4\% of the gold spans), type confusion is by far the
most common DistilBERT error. Inspection of the recorded examples
reveals classic CoNLL ambiguities: country names used metonymically for
national teams (\textsc{LOC} tagged as \textsc{ORG}), surnames identical
to place names (\textsc{LOC} vs.\ \textsc{PER}), and ambiguous
sentence-initial proper nouns. These errors require world knowledge
that cannot be extracted from the 14k-sentence training set alone.

\paragraph{False positives concentrate in \textsc{MISC}.}
All three of the recorded false-positive examples are \textsc{MISC}
predictions over tokens that are not entities at all. \textsc{MISC} is a
catch-all class for nationality adjectives, event names, and other
``odd'' proper nouns; the model has learnt to over-predict it whenever
capitalisation cues are present without strong contextual constraints.

\subsection{Discussion: Role of Contextual Embeddings}

The BiLSTM-CRF must learn word representations from scratch using only
the 14k-sentence training corpus. Out-of-vocabulary words are mapped to
a single \texttt{<UNK>} embedding, which severely limits generalisation
to rare or unseen proper nouns --- precisely the population that NER
must handle well. Even with a structured CRF decoder enforcing label
consistency, the model achieves only $0.7011$ F1.

DistilBERT's contextual sub-word embeddings address both limitations
simultaneously. WordPiece tokenisation eliminates the closed-vocabulary
problem entirely: any character sequence is decomposable into known
sub-units, allowing the model to construct meaningful representations
for \textit{Schmidt}, \textit{Kyiv}, or \textit{Chrysler} without ever
having seen those tokens during pre-training. Self-attention provides
global context for every position, so the same surface form can receive
different representations depending on its context (e.g.\
\textit{Washington} as a \textsc{PER} versus a \textsc{LOC}). These two
properties combined explain the almost 20-point F1 gap between the two
models.

The principal trade-off is computational: DistilBERT requires roughly an
order of magnitude more memory and compute than the BiLSTM-CRF for both
training and inference, even in its distilled form. Where the latency
or memory budget is tight, BiLSTM-CRF remains a reasonable default;
where accuracy is paramount, contextual representations are clearly
superior.
\newpage
% ─────────────────────────────────────────────────────────────────
% CENG 467 Midterm – Question 3: Text Summarization
% Student ID: 310201050
% ─────────────────────────────────────────────────────────────────

\section{Question 3: Text Summarization}

\subsection{Dataset}

The CNN/DailyMail v3.0.0 corpus \cite{see2017get} is a large-scale
news-summarisation benchmark consisting of articles from CNN and the
Daily Mail paired with multi-sentence editor-written highlights. The
full corpus contains $\sim$287k training, 13k validation, and 11k test
articles. For this study, the first $200$ examples of the test split
are used (selected deterministically with seed 3102). The mean article
length in the subset is $555$ words; the mean reference summary length
is $34$ words across roughly 3 sentences. These characteristics motivate
setting LexRank's sentence budget to 3.

\subsection{Approaches}

\subsubsection{LexRank (Extractive, Graph-Based)}

LexRank \cite{erkan2004lexrank} represents each sentence in an article as
a node in an undirected graph, with edge weights equal to TF-IDF cosine
similarity between sentence pairs. A PageRank-style stationary
distribution is then computed, and the $k$ highest-ranking sentences are
returned verbatim as the summary. Because LexRank operates by selecting
existing sentences, it cannot introduce unfaithful content (every output
token comes from the source document), but it also cannot rephrase,
compress, or fuse information across sentences.

We use the implementation provided by the \texttt{sumy} library with the
English Snowball stemmer and standard English stopwords. Three sentences
are extracted per article, matching the average sentence count of the
gold-standard highlights.

\subsubsection{DistilBART (Abstractive, Encoder–Decoder Transformer)}

We use \texttt{sshleifer/distilbart-cnn-12-6}, a distilled BART
\cite{lewis2020bart} checkpoint with 12 encoder layers and 6 decoder
layers that has been pre-fine-tuned on CNN/DailyMail. The choice of
distilled weights is dictated by the M2 8\,GB memory budget. Generation
hyperparameters are summarised in Table~\ref{tab:bart_gen}.

\begin{table}[h]
\centering
\caption{DistilBART Generation Configuration}
\label{tab:bart_gen}
\begin{tabular}{ll}
\toprule
Parameter & Value \\
\midrule
Maximum input length        & 1024 sub-tokens (truncation) \\
Minimum / maximum output len & 30 / 130 sub-tokens \\
Decoding strategy           & Beam search, $B = 4$ \\
Length penalty              & 2.0 \\
No-repeat $n$-gram size     & 3 \\
Early stopping              & enabled \\
\bottomrule
\end{tabular}
\end{table}

\subsection{Evaluation Metrics}

Summaries are evaluated against the gold highlights using a multi-faceted
suite that captures complementary aspects of quality:

\begin{itemize}
  \item \textbf{ROUGE-1, ROUGE-2, ROUGE-L} \cite{lin2004rouge}: $F_1$ of
        unigram, bigram, and longest-common-subsequence overlap.
        Standard for summarisation but rewards lexical overlap rather
        than meaning.
  \item \textbf{BLEU} \cite{papineni2002bleu}: $n$-gram precision with
        brevity penalty (smoothing method 4). Borrowed from machine
        translation; tends to penalise abstractive paraphrasing.
  \item \textbf{METEOR} \cite{banerjee2005meteor}: incorporates stemming
        and synonym matching, partially correcting for ROUGE/BLEU's
        surface-level rigidity.
  \item \textbf{BERTScore F1} \cite{zhang2020bertscore}: token-level
        cosine similarity in contextual embedding space
        (\texttt{distilbert-base-uncased} backbone for memory
        efficiency). Captures semantic equivalence even under heavy
        paraphrasing.
\end{itemize}

\subsection{Results}

\begin{table}[h]
\centering
\caption{Summarisation Results on CNN/DailyMail (n=200 test articles)}
\label{tab:q3_results}
\begin{tabular}{lcc}
\toprule
Metric              & LexRank (extractive) & DistilBART (abstractive) \\
\midrule
ROUGE-1             & 0.2808 & \textbf{0.3538} \\
ROUGE-2             & 0.0939 & \textbf{0.1498} \\
ROUGE-L             & 0.1883 & \textbf{0.2580} \\
BLEU                & 0.0553 & \textbf{0.1043} \\
METEOR              & 0.3075 & \textbf{0.3370} \\
BERTScore F1        & 0.7776 & \textbf{0.8030} \\
\bottomrule
\end{tabular}
\end{table}

DistilBART outperforms LexRank on every single metric. The gap is
largest for $n$-gram-overlap metrics that reward lexical fidelity to the
reference (ROUGE-2 nearly $1.6\times$ higher; BLEU nearly $2\times$
higher). The gap on METEOR and BERTScore is narrower because both
metrics partially reward LexRank's verbatim copying of source-document
phrasings.

\subsection{Qualitative Analysis}

Three example articles are inspected in full in
\texttt{q3\_results.json} (indices 0, 50, 100). Three differences
between the extractive and abstractive outputs are consistently observed.

\paragraph{Fluency.}
LexRank summaries are by construction grammatical sentences from the
original article, so they are individually fluent. However, when the
three top-ranked sentences come from non-adjacent positions in the
source, the resulting concatenation reads as disjointed narrative ---
the test-set example at index 50 illustrates this: LexRank concatenates
three legal-discussion sentences about defamation that lose any
connection to the article's main narrative about the Rolling Stone
retraction. DistilBART, in contrast, produces a coherent paragraph
because the decoder generates it autoregressively under a smooth
language-model prior.

\paragraph{Information Coverage.}
DistilBART recovers the article's most newsworthy facts (\textit{``A new
trial was ordered in 2014 after firearms experts testified''} in example
100) that LexRank misses entirely, because such facts are spread across
multiple low-salience sentences in the source. LexRank's centrality-based
selection rewards repeated themes rather than rare informative content.

\paragraph{Factual Faithfulness.}
This is where the trade-off reverses. LexRank's outputs are extractively
faithful by definition. DistilBART, in contrast, occasionally exhibits
\emph{hallucination}: in example 50, the abstractive output attributes
the article to a named columnist (``Dean Obeidallah'') whose name does
not appear in the source preview. Such hallucinations are well
documented for abstractive encoder–decoder models trained with
maximum-likelihood objectives, and they remain a key open problem for
the field.

\subsection{Discussion}

ROUGE and BLEU favour systems that lexically overlap with the reference,
so LexRank --- which copies source sentences directly --- typically
scores respectably on these metrics despite generating less coherent
prose. However, the present results show that even on these surface
metrics, abstractive DistilBART wins decisively, because its
fine-tuning on CNN/DailyMail-style highlights teaches it to produce
output that closely matches the lexical distribution of human-written
summaries. DistilBART's lead widens further on METEOR and BERTScore,
both of which reward semantic alignment beyond surface form.

In terms of cost, LexRank requires only TF-IDF computation and a small
eigenvector calculation, runs entirely on CPU, and produces summaries
in milliseconds. DistilBART requires a fine-tuned $\sim$300\,MB model
and beam-search decoding that takes seconds per article. On readability
and coverage, DistilBART is clearly preferable. On faithfulness,
LexRank is strictly safer.

In production settings, the choice depends heavily on downstream
tolerance for hallucination: legal, medical, and financial
summarisation favour extractive (or strongly grounded abstractive)
approaches, whereas consumer-facing news digests can usually tolerate
the abstractive trade-off in exchange for fluency.
\newpage
% ─────────────────────────────────────────────────────────────────
% CENG 467 Midterm – Question 4: Machine Translation
% Student ID: 310201050
% ─────────────────────────────────────────────────────────────────

\section{Question 4: Machine Translation}

\subsection{Dataset and Preprocessing}

The Multi30k corpus \cite{elliott2016multi30k} is used for
German-to-English translation. Multi30k is built on top of the
Flickr30k image-caption dataset, with each English caption
professionally translated into German. It contains 29{,}000 training,
1{,}014 validation, and 1{,}000 test parallel sentence pairs. The
relatively short, syntactically-controlled captions make the corpus a
popular benchmark for neural-MT experiments on modest hardware.

Two distinct preprocessing pipelines are used:

\begin{itemize}
  \item \textbf{Seq2Seq pipeline:} text is lowercased, punctuation
        tokens are separated from adjacent words via regular
        expressions, and resulting tokens are mapped to integer ids
        using vocabularies built from the training data. Vocabulary
        size is capped at 8{,}000 per language with a minimum-frequency
        threshold of 2; rare tokens are mapped to \texttt{<unk>}. After
        capping, the actual learnt vocabularies contain
        $7{,}835$ German source tokens and $5{,}895$ English target
        tokens. Sequences are padded to a maximum of 50 tokens, and
        decoder targets are wrapped in \texttt{<sos>}/\texttt{<eos>}
        markers.
  \item \textbf{MarianMT pipeline:} the pretrained
        \texttt{Helsinki-NLP/opus-mt-de-en} tokenizer
        (SentencePiece-based) is applied directly, requiring no
        additional preprocessing.
\end{itemize}

\subsection{Models}

\subsubsection{Model A — Seq2Seq with Bahdanau Attention (Trained from Scratch)}

A bidirectional LSTM encoder maps source tokens to a context tensor; the
decoder is a unidirectional LSTM with Bahdanau attention
\cite{bahdanau2015neural}, where the attention layer computes a soft
alignment $\alpha_t$ over encoder states at every decoding step:

\[
\alpha_t = \mathrm{softmax}\bigl(v^{\top}\tanh(W[h_t^{\text{dec}};h_{1:T}^{\text{enc}}])\bigr),
\quad
c_t = \sum_{i=1}^{T}\alpha_{t,i}\, h_i^{\text{enc}}.
\]

The output projection conditions on the LSTM state, the attention
context, and the decoder embedding, following the standard formulation.
Hyperparameters are summarised in Table~\ref{tab:s2s_hp}; the trained
model has $12.12$ million parameters.

\begin{table}[h]
\centering
\caption{Seq2Seq+Attention Hyperparameters}
\label{tab:s2s_hp}
\begin{tabular}{ll}
\toprule
Hyperparameter      & Value \\
\midrule
Embedding dimension & 256 \\
Hidden dimension    & 256 (encoder bidirectional → 512 effective) \\
Dropout             & 0.30 \\
Batch size          & 32 \\
Epochs              & 8 \\
Optimiser / LR      & Adam / $10^{-3}$, StepLR (step 3, $\gamma=0.5$) \\
Teacher-forcing     & 0.50 \\
Gradient clipping   & 1.0 \\
Decoding            & Greedy \\
Total parameters    & 12.12 M \\
\bottomrule
\end{tabular}
\end{table}

\subsubsection{Model B — MarianMT (Pretrained Transformer)}

The \texttt{Helsinki-NLP/opus-mt-de-en} checkpoint
\cite{tiedemann2020opusmt} is a 6-layer encoder–decoder Transformer
trained on the OPUS corpora across many language pairs. The model is
used out-of-the-box for inference (no additional fine-tuning); generation
uses beam search with $B=4$ and a maximum output length of 128
sub-tokens.

\subsection{Evaluation Metrics}

We evaluate translations against the gold English captions using a
four-metric suite chosen to capture complementary aspects of quality:

\begin{itemize}
  \item \textbf{BLEU} \cite{papineni2002bleu} via \texttt{sacrebleu}.
        Standard $n$-gram precision with brevity penalty; reproducible
        exactly across toolchains.
  \item \textbf{METEOR} \cite{banerjee2005meteor}. Aligns hypotheses to
        references using stemming and WordNet synonyms.
  \item \textbf{ChrF} \cite{popovic2015chrf} via \texttt{sacrebleu}.
        Character $n$-gram F-score; particularly informative for
        morphologically rich languages because it credits stem matches
        even when affixes differ.
  \item \textbf{BERTScore F1} \cite{zhang2020bertscore} computed with a
        \texttt{distilbert-base-uncased} backbone (for memory
        efficiency). Captures semantic equivalence in contextual
        embedding space.
\end{itemize}

\subsection{Results}

\begin{table}[h]
\centering
\caption{Multi30k Test-Set Translation Results (de→en, n=1000)}
\label{tab:q4_results}
\begin{tabular}{lcc}
\toprule
Metric            & Seq2Seq+Attention & MarianMT \\
\midrule
BLEU              & 27.69 & \textbf{39.66} \\
ChrF              & 49.34 & \textbf{60.60} \\
METEOR            & 0.6267 & \textbf{0.7394} \\
BERTScore F1      & 0.8882 & \textbf{0.9333} \\
\bottomrule
\end{tabular}
\end{table}

MarianMT outperforms the from-scratch Seq2Seq across every metric. The
margin is particularly pronounced on BLEU ($+11.97$ points) and ChrF
($+11.26$ points), which reward exact $n$-gram matches and are
therefore especially sensitive to fluency and lexical accuracy. The
narrower margin on BERTScore (still $+0.045$) indicates that Seq2Seq
captures most of the semantic content but expresses it less fluently.

\subsection{Qualitative Examples}

\begin{table}[h]
\centering
\caption{Sample Outputs (test indices 0, 100, 500)}
\label{tab:q4_qual}
\small
\begin{tabular}{p{0.16\textwidth}p{0.74\textwidth}}
\toprule
\multicolumn{2}{l}{\textbf{Example 1 (test index 0)}} \\
\midrule
Source (de)  & Ein Mann mit einem orangefarbenen Hut, der etwas anstarrt. \\
Reference    & A man in an orange hat starring at something. \\
Seq2Seq+Attn & a man in an orange hat is something something . \\
MarianMT     & A man with an orange hat staring at something. \\
\midrule
\multicolumn{2}{l}{\textbf{Example 2 (test index 100)}} \\
\midrule
Source (de)  & Eine glückliche Frau bereitet in einem Coffee-Shop eine Erfrischung zu. \\
Reference    & A happy woman is preparing a refreshment at a coffee shop. \\
Seq2Seq+Attn & a happy is preparing a a a a a a a a . \\
MarianMT     & A happy woman prepares a refreshment in a coffee shop. \\
\midrule
\multicolumn{2}{l}{\textbf{Example 3 (test index 500)}} \\
\midrule
Source (de)  & Ein Tisch voller gerahmter Bilder, ein Freiluftmarkt. \\
Reference    & A table full of pictures in frames, at an outdoor market. \\
Seq2Seq+Attn & a table full of pictures pictures . \\
MarianMT     & A table full of framed pictures, an outdoor market. \\
\bottomrule
\end{tabular}
\end{table}

Across these three samples, three patterns are consistently observed.

\paragraph{Repetition collapse in Seq2Seq.}
Example 100 reveals a classic failure mode of greedy decoding in small
recurrent models: the model gets stuck on the article \textit{a} and
emits it eight times in succession. This is the well-documented
``repetition trap'' that motivated the development of coverage
mechanisms and beam search with length-penalty constraints. The same
phenomenon appears in Example 1 (\textit{``something something''}) and
Example 3 (\textit{``pictures pictures''}).

\paragraph{Vocabulary clipping artefacts.}
Example 100 also shows the Seq2Seq model dropping the noun
\textit{Frau} (woman): the German source word likely received a
sufficient embedding, but the target-side decoder failed to reproduce
\textit{woman} in this context. MarianMT's SentencePiece tokenizer and
larger pre-training corpus produce \textit{woman} fluently.

\paragraph{Surface fluency.}
Even when the Seq2Seq output is semantically close, it can produce
slightly ungrammatical English (article omission, lowercased proper
sentence starts). MarianMT outputs are uniformly fluent because the
decoder language model has been trained on orders of magnitude more
data.

\subsection{Discussion}

\paragraph{What each metric reflects.}
BLEU and ChrF reward surface-level overlap, with BLEU penalising
paraphrases harshly and ChrF being more forgiving thanks to
character-level matching --- particularly relevant for German
morphology. METEOR partially recovers paraphrases through stemming and
synonymy, and tends to track human judgement more closely than BLEU
alone. BERTScore is the most permissive: it credits semantic
equivalence even when surface forms differ substantially, which is why
the Seq2Seq model retains a relatively high BERTScore ($0.8882$)
despite its low BLEU ($27.69$).

\paragraph{Architectural trade-offs.}
The Seq2Seq+attention model has $12.12$M parameters, trains in roughly
ten minutes on M2 hardware, and runs inference in seconds. Its weakness
is data efficiency: with only 29k training pairs and randomly
initialised embeddings, it cannot match the quality of a model that has
seen orders of magnitude more bilingual data. MarianMT inherits roughly
$78$M parameters of pretrained knowledge, handles open-vocabulary
inputs through SentencePiece, and uses self-attention rather than
recurrence, all of which make it both more accurate and more robust.

The trade-off is in resource use and architectural transparency: a
hand-written Seq2Seq+attention model is a valuable pedagogical artefact
precisely because every component is visible and modifiable, whereas
the pretrained Marian checkpoint is a black box that one consumes
rather than understands.
\newpage
% ─────────────────────────────────────────────────────────────────
% CENG 467 Midterm – Question 5: Language Modeling
% Student ID: 310201050
% ─────────────────────────────────────────────────────────────────

\section{Question 5: Language Modeling}

\subsection{Dataset}

The WikiText-2 corpus \cite{merity2017wikitext} is selected for this
investigation. WikiText-2 contains approximately 2 million tokens drawn
from a curated set of verified-good Wikipedia articles. The corpus is
preprocessed into fixed train ($2{,}051{,}910$ tokens), validation
($213{,}886$ tokens) and test ($241{,}211$ tokens) splits, and
words occurring fewer than three times in the training data have been
replaced with the \texttt{<unk>} symbol. These properties make
WikiText-2 the canonical small benchmark for language modelling
experiments and a natural choice for an Apple~M2 8\,GB development
setup.

The vocabulary is constructed directly from the training tokens
($33{,}277$ unique types) and shared by both models, ensuring that
perplexity numbers are directly comparable.

\subsection{Models}

\subsubsection{Model 1 — Trigram with Add-k Smoothing}

The first model is a classical trigram language model with add-$k$
smoothing. Letting $c(\cdot)$ denote training-corpus counts, the
probability of token $w_3$ given two-word context $(w_1, w_2)$ is

\[
P_{\text{add-}k}(w_3 \mid w_1, w_2) =
\frac{c(w_1, w_2, w_3) + k}{c(w_1, w_2) + k\,V},
\]

where $V$ is the vocabulary size and $k = 0.01$. This small $k$ assigns
non-zero probability to unseen trigrams while disturbing the
distribution over observed contexts only slightly. Compared with
classical add-1 (Laplace) smoothing it preserves more of the empirical
signal in high-count contexts.

The model has no learned parameters in the usual sense; estimation is a
single pass over the training corpus. Generation is performed by
sampling from the smoothed distribution at each step.

\subsubsection{Model 2 — 2-Layer LSTM Language Model}

The second model is a two-layer LSTM \cite{hochreiter1997lstm} trained
end-to-end with cross-entropy loss using truncated back-propagation
through time (BPTT). The trained model has $15.93$ million parameters.

\begin{table}[h]
\centering
\caption{LSTM Language Model Hyperparameters}
\label{tab:lstm_lm_hp}
\begin{tabular}{ll}
\toprule
Hyperparameter   & Value \\
\midrule
Embedding dimension      & 192 \\
Hidden dimension         & 256 \\
Number of LSTM layers    & 2 \\
Dropout                  & 0.45 \\
BPTT context length      & 35 tokens \\
Batch size               & 20 (BPTT layout) \\
Epochs                   & 6 \\
Optimiser / LR           & Adam / $10^{-3}$ \\
LR schedule              & StepLR (step 2, $\gamma=0.5$) \\
Gradient clipping        & 0.25 \\
Total parameters         & 15.93 M \\
\bottomrule
\end{tabular}
\end{table}

The hidden state is detached at every BPTT window so gradients do not
back-propagate across windows, which both stabilises training and keeps
memory consumption bounded. At inference, a fresh zero hidden state
followed by the prompt tokens primes the model before sampling begins.

\subsection{Evaluation: Perplexity}

Perplexity is the exponentiated per-token cross-entropy:

\[
\mathrm{PPL}(w_1,\dots,w_N) = \exp\Bigl(-\tfrac{1}{N}\sum_{i=1}^{N} \log P(w_i \mid w_{<i})\Bigr).
\]

It is the standard intrinsic measure of language-model quality; lower
values indicate that the model assigns higher probability mass to
held-out text.

\begin{table}[h]
\centering
\caption{WikiText-2 Perplexity (lower is better; seed 3102)}
\label{tab:q5_results}
\begin{tabular}{lcc}
\toprule
Model                            & Validation PPL & Test PPL \\
\midrule
Trigram (add-$k$, $k = 0.01$)    & 3868.40 & 3587.97 \\
2-layer LSTM                     & \textbf{389.69} & \textbf{353.41} \\
\bottomrule
\end{tabular}
\end{table}

The LSTM achieves perplexity approximately \emph{ten times lower} than
the trigram baseline ($353.41$ vs.\ $3587.97$ on test). This dramatic
gap quantifies the benefit of distributed representations and longer
effective context for next-token prediction.

\subsection{Text Generation Samples}

Both models are seeded with three identical short prompts and generate
25 follow-up tokens by sampling from their respective distributions
(temperature $\tau = 0.9$ for the LSTM; raw smoothed distribution for the
trigram).

\begin{table}[h]
\centering
\caption{Generated Continuations (excerpts)}
\label{tab:q5_gen}
\small
\begin{tabular}{p{0.13\textwidth}p{0.78\textwidth}}
\toprule
Prompt    & ``the city of'' \\
Trigram   & \textit{... Zagreb. However, German mycologist Rolf Singer listed three subspecies, although it was then tied with Tower of Terror, with Julio Jones} \\
LSTM      & \textit{... achieving to October the kings before the soul. There are full, was \texttt{<unk>}, almost serving as the Bruins as a short school} \\
\midrule
Prompt    & ``in the year'' \\
Trigram   & \textit{... . Twenty @-@ four @-@ issue story arc to be in a Pepsi television advertisement in the north of the young characters into one of} \\
LSTM      & \textit{... – two season and the 17th of arguing and fought on Scotland and ₹ 85 (1 @.@) mi; the storm – 16} \\
\midrule
Prompt    & ``she said that'' \\
Trigram   & \textit{... the problem in many cases from several teams that were arrested and killed four young black girls and football manager. Lawton and Sewell's} \\
LSTM      & \textit{... it is `` that related is while deal as his political titles '', by the first name on six side, being peaking by} \\
\bottomrule
\end{tabular}
\end{table}

\paragraph{Trigram outputs.}
Locally fluent two- and three-word fragments appear because every chosen
word has been observed immediately following the preceding bigram.
However, coherence collapses beyond the trigram window: the topic
drifts unpredictably across very different domains within a single
sample (e.g.\ from \textit{German mycologist} to \textit{Tower of
Terror} to \textit{Julio Jones}), and named entities are introduced
without context. Topical coherence is bounded by the length of the
$n$-gram context and breaks immediately at that boundary.

\paragraph{LSTM outputs.}
The LSTM produces longer fluent spans with consistent grammatical
structure (\textit{``almost serving as the Bruins as a short school''}
preserves clausal structure even where content is incoherent). Failure
modes are different from the trigram model's: rather than topic drift,
the LSTM tends to combine plausible noun phrases in semantically
implausible ways, and occasionally produces tokens such as \texttt{₹},
\texttt{@-@}, and \texttt{<unk>} that are artefacts of the WikiText-2
preprocessing rather than genuine language. These behaviours are
characteristic of LM-decoded text trained with maximum-likelihood
objectives on a small corpus.

\subsection{Discussion}

\paragraph{Why the LSTM dominates on perplexity.}
The trigram model estimates one parameter per observed $(w_1, w_2,
w_3)$ triple, which suffers from severe data sparsity: even on
WikiText-2's $\sim$2M training tokens, the vast majority of
theoretically possible trigrams remain unobserved, so add-$k$ smoothing
must allocate substantial probability mass to all of them uniformly.
The LSTM, in contrast, learns a continuous distributed representation
in which similar contexts produce similar predictions, generalising
automatically across rare or unseen surface patterns. The
$10\times$ perplexity gap quantifies this advantage.

\paragraph{Effective context.}
Although the LSTM is in principle able to condition on arbitrarily long
histories, in practice the information that survives many recurrent
updates degrades. Empirical studies have shown effective context
lengths of a few hundred tokens for deep LSTMs; modern Transformers use
direct attention to circumvent this limitation, which is the main
reason they have largely supplanted recurrent architectures for
language modelling. The trigram model's effective context is fixed at
two tokens, which is why its perplexity remains in the four-digit range
even after smoothing.

\paragraph{Trade-offs.}
The trigram model trains in seconds, requires no GPU, and is perfectly
transparent --- every probability is a counted fraction. The LSTM
requires multiple training epochs and approximately $16$M parameters,
but produces dramatically lower perplexity and generates more coherent
samples. For applications that need quick statistical baselines (corpus
statistics, code-switch detection, forced alignment) a smoothed
$n$-gram remains a perfectly serviceable choice; for any application
that requires generation quality, the neural model is preferable.

% =============================================================================
% CONCLUSION
% =============================================================================
\newpage
\section{Conclusion}

Across the five tasks studied, a consistent pattern emerges:
contextual sub-word representations from pretrained Transformers
outperform both classical statistical baselines and recurrent
architectures trained from scratch, often by a substantial margin.
The advantage is most pronounced when training data is limited
(NER, MT) and when surface lexical features are weakly informative
(language modelling). On the IMDb classification task, where the
6{,}000-example training subset is small but the unigram signal is
strong, the sparse TF-IDF baseline came surprisingly close to
DistilBERT, even outperforming the from-scratch BiLSTM.

The study also highlights three trade-offs that should temper any
universal endorsement of contextual representations. First,
inference cost: pretrained Transformers are dramatically more
expensive than $n$-gram or LSTM baselines, a fact that became
concrete during these experiments when the DistilBERT fine-tune for
Q1 alone consumed roughly 45\,minutes on an M2. Second,
faithfulness: abstractive Transformer decoders can hallucinate
content not present in the source, a failure mode that extractive
baselines structurally cannot exhibit (Q3, example 50). Third,
transparency: a from-scratch model whose parameters one has
trained, debugged, and inspected is qualitatively different from a
downloaded checkpoint, and the educational and diagnostic value of
the former is non-trivial. The Q4 Seq2Seq model exhibited a
classic ``repetition trap'' (\textit{``a a a a a''}) in three of
three qualitative samples, a failure mode that is invisible when
one only consumes a pretrained black box.

Reproducibility was treated as a first-class constraint throughout.
Every result in this report can be regenerated by running the
provided scripts on a comparable device with seed 3102. The
associated repository contains the JSON output files used to
populate every numeric cell in this report.

% =============================================================================
% REFERENCES
% =============================================================================
\newpage
\bibliographystyle{plainnat}
\bibliography{references}

\end{document}